%------------------------------------------------------------------------------%
\documentclass[a4paper,12pt,fleqn]{article}

\usepackage{calculo_varias_variaveis-1}

\ShowAnswers

%------------------------------------------------------------------------------%
\begin{document}

\makeHeader{CFVVI}{Prova 1 -- Turma 8}{01/04/2026}

% 01
%------------------------------------------------------------------------------%
\exercise{50\%}
Uma partícula descreve uma trajetória senoidal descrita
pelas equações paramétricas
\[
  x(t) = 6 - 3t
  \qquad\qquad
  y(t) = 4\cos\left(\frac{\pi t}{2}\right) + 2
  \qquad\qquad\qquad
  0 \leq t \leq 2
\]
\begin{enumerate}[label=\alph*),itemsep=0pt,topsep=0pt]
  \item\!\! [10\%] determine on pontos onde da curva associados com
  $t=0$, $t=1$ e em $t=2$
  \item\!\! [15\%] calcule a expressão para o vetor tangente à curva
  \item\!\! [15\%] avalie o vetor tangente em
  $t=0$, $t=1$ e em $t=2$
  \item\!\! [10\%] esboce a curva e os vetores
  tangentes calculados
  (considere $\pi\approx 3$)
\end{enumerate}
~

\begin{questiononly}
  \qquad\includegraphics{axis_xy}
\end{questiononly}
\begin{answeronly}
  \qquad\includegraphics{T8_B_parametrica}
\end{answeronly}
\clearpagequestiononly

\begin{answer}
  \textbf{a)} \quad
  Calculando os pontos associados a $t=0$, $t=1$ e em $t=2$
  \begin{multicols}{2}
    \begin{align*}
      x(0)
      & = \left(6 - 3t\right)\evalat{0}{} \\
      & = 6 - 3\times 0 \\
      & = 6
    \end{align*}
    \newcolumn
    \begin{align*}
      y(0)
      & = \left(4\cos\left(\frac{\pi t}{2}\right) + 2 \right)\evalat{0}{} \\
      & = 4\cos\left(\frac{\pi \times 0}{2}\right) + 2 \\
      & = 4 \times 1 + 2 \\
      & = 6
    \end{align*}
  \end{multicols}
  \begin{multicols}{2}
    \begin{align*}
      x(1)
      & = \left(6 - 3t\right)\evalat{1}{} \\
      & = 6 - 3\times 1 \\
      & = 3
    \end{align*}
    \newcolumn
    \begin{align*}
      y(1)
      & = \left(4\cos\left(\frac{\pi t}{2}\right) + 2 \right)\evalat{1}{} \\
      & = 4\cos\left(\frac{\pi}{2}\right) + 2 \\
      & = 4 \times 0 + 2 \\
      & = 2
    \end{align*}
  \end{multicols}
  \begin{multicols}{2}
    \begin{align*}
      x(2)
      & = \left(6 - 3t\right)\evalat{2}{} \\
      & = 6 - 3\times 2 \\
      & = 0
    \end{align*}
    \newcolumn
    \begin{align*}
      y(2)
      & = \left(4\cos\left(\frac{\pi t}{2}\right) + 2 \right)\evalat{2}{} \\
      & = 4\cos\left(\frac{\pi \times 2}{2}\right) + 2 \\
      & = 4(-1) + 2 \\
      & = -2
    \end{align*}
  \end{multicols}

  \vspace{-\baselineskip}
  \textbf{b)} \quad
  O vetor tangente é
  \(
  v(t) = \Vetor{\frac{dx}{dt}}{\frac{dy}{dt}}
  \)

  Avaliando cada derivada
  \begin{multicols}{2}
    \[
      \frac{dx}{dt}
      = \frac{d}{dt}\left[6 - 3t\right]
      = -3
    \]
    \newcolumn
    \begin{align*}
      \frac{dy}{dt}
      & = \frac{d}{dt}\left[4\cos\left(\frac{\pi t}{2}\right) + 2\right] \\
      & = -4\sin\left(\frac{\pi t}{2}\right)\frac{d}{dt}\left(\frac{\pi t}{2}\right) \\
      & = -4\sin\left(\frac{\pi t}{2}\right)\frac{\pi}{2} \\
      & = -2\pi\sin\left(\frac{\pi t}{2}\right)
    \end{align*}
  \end{multicols}
  \vspace{-\baselineskip}
  Portanto,
  \(
    v(t) = \vetor{-3}{-2\pi\sin\left(\frac{\pi t}{2}\right)}
  \)

  \bigskip
  \textbf{c)} \quad
  Avaliando o vetor tangente nos pontos $t=0$, $t=1$ e em $t=2$
  \begin{multicols}{3}
    \setlength{\mathindent}{0pt}
    \begin{align*}
      v(0)
      & = \Vetor{-3}{-2\pi\sin\left(\frac{\pi t}{2}\right)} \evalat{0}{} \\[2mm]
      & = \Vetor{-3}{-2\pi\sin\left(0\right)} \\[2mm]
      & = \vetor{-3}{-2\pi\times 0} \\[2mm]
      & = \vetor{-3}{0}
    \end{align*}

    \newcolumn
    \begin{align*}
      v(1)
      & = \Vetor{-3}{-2\pi\sin\left(\frac{\pi t}{2}\right)} \evalat{1}{} \\[2mm]
      & = \Vetor{-3}{-2\pi\sin\left(\frac{\pi}{2}\right)} \\[2mm]
      & = \Vetor{-3}{-2\pi\times 1} \\[2mm]
      & = \vetor{-3}{-2\pi}
    \end{align*}

    \newcolumn
    \begin{align*}
      v(2)
      & = \Vetor{-3}{-2\pi\sin\left(\frac{\pi t}{2}\right)} \evalat{2}{} \\[2mm]
      & = \Vetor{-3}{-2\pi\sin\left(\pi\right)} \\[2mm]
      & = \Vetor{-3}{-2\pi\times 0} \\[2mm]
      & = \vetor{-3}{0}
    \end{align*}
  \end{multicols}
\end{answer}


% 02
%------------------------------------------------------------------------------%
\exercise{50\%}
Considerando a função
\(
  f(x, y) = \sqrt{\frac{2}{9 - y^2}}
\)
\begin{enumerate}[label=(\alph*),itemsep=0pt]
  \item \ [15\%] determine o domínio de $f$
  \item \ [15\%] esboce o domínio
  \item \ [10\%] determine a curva de nível que passa pelo ponto \((2, 2)\)
  \item \ [10\%] esboce a curva de nível que passa pelo ponto \((2, 2)\)
\end{enumerate}
\begin{questiononly}
  \begin{multicols}{2}
    \;\; Domínio
    \vspace{-0.3\baselineskip}
    \begin{center}
      \includegraphics{axis_xy}
    \end{center}
    \;\; Curva de nível
    \vspace{-0.3\baselineskip}
    \begin{center}
      \includegraphics{axis_xy}
    \end{center}
  \end{multicols}
\end{questiononly}
\begin{answeronly}
  \begin{multicols}{2}
    \;\; Domínio
    \vspace{-0.3\baselineskip}
    \begin{center}
      \includegraphics{T8_B_dominio}
    \end{center}
    \;\; Curva de nível
    \vspace{-0.3\baselineskip}
    \begin{center}
      \includegraphics{T8_B_curva_nivel}
    \end{center}
  \end{multicols}
\end{answeronly}

\begin{answer}
  \textbf{a)} \quad
  Não podemos calcular a raiz quadrada de um número negativo e não
  podemos dividir por zero, portanto, o domínio de $f$ é o conjunto
  de todos os pontos \((x, y)\in\R^2\) tais que
  \[
    \frac{2}{9 - y^2} \geq 0 \qquad\text{ e }\qquad 9 - y^2 \text{ não é zero}
  \]
  portanto
  \begin{align*}
    9 - y^2 & > 0            \\
    -y^2    & > -9           \\
    y^2     & < 9            \\
    \abs{y} & < \sqrt{9} = 3 \\
    -3 < y  & < 3
  \end{align*}
  ou seja, os pontos entre as retas horizontais \(y=-3\) e \(y=3\).
  Assim, o domínio é
  \[
    D_f = \left\{(x, y) \in \R^2 \st -3 < y < 3 \right\}
  \]

  \bigskip
  \textbf{c)} \quad
  A curva de nível que passa pelo ponto $(2, 2)$ tem o valor
  \[
    c
    = f(2, 2)
    = \sqrt{\frac{2}{9 - y^2}} \evalat{(2, 2)}{}
    = \sqrt{\frac{2}{9 - 2^2}}
    = \sqrt{\frac{2}{5}}
  \]
  A curva de nível associada é
  \begin{align*}
    f(x, y)                  & = \sqrt{\frac{2}{5}} \\[2mm]
    \sqrt{\frac{2}{9 - y^2}} & = \sqrt{\frac{2}{5}} \\[2mm]
    \frac{2}{9 - y^2}        & = \frac{2}{5}        \\[2mm]
    9 - y^2                  & = 5                  \\[2mm]
    - y^2                    & = 5 - 9 = -4         \\[2mm]
    y^2                      & = 4                  \\[2mm]
    \abs{y}                  & = 2                  \\[2mm]
    y                        & = \pm 2
  \end{align*}
  A curva de nível é o par de retas horizontais
  \[
    y = 2
    \qquad\text{ e }\qquad
    y = -2
  \]

\end{answer}

%------------------------------------------------------------------------------%
\end{document}
%------------------------------------------------------------------------------%
